%\ifodd\value{page}\hbox{}\newpage\fi
\chapter*{Śrī Lalitā Sahasranāma — Śloka 9}
\addcontentsline{toc}{chapter}{Śloka 9 - Nomi 20,21}

\begin{sloka}
{\sanskritfont
\begin{tabular}{c}
पद्मरागशिलादर्शपरिभाविकपोलभूः ।\\
नवविद्रुमबिम्बश्रीन्यक्कारिरदनच्छदा ॥ ९॥
\end{tabular}

\vspace{1.5em}

{\middlesize \iastfont
padmarāga-śilā-darśa-paribhāvi-kapola-bhūḥ |\\
nava-vidruma-bimba-śrī-nyakkāri-radana-cchadā ||9||}

\vspace{1em}

{\middlesize
\anvaya{%
पद्मरागशिलादर्शपरिभाविकपोलभूः ।
नवविद्रुमबिम्बश्रीन्यक्कारिरदनच्छदा ।
}
}

\middlesize \iastfont \itmean{%
«Le cui guance mettono in ombra lo splendore riflettente della pietra di \textit{padmarāga};  
il cui rivestimento dei denti fa sfigurare la bellezza del corallo fresco e del frutto \textit{bimba}.»%
}

}
\end{sloka}

\wordmeaning{%
\wordline{पद्मराग-शिला-दर्श-परिभावि-कपोल-भूः}{padmarāga-śilā-darśa-paribhāvi-\\ kapola-bhūḥ}
  { le cui guance (\textit{kapola}) sono tali da mettere in ombra (\textit{paribhāvi}) lo splendore riflettente (\textit{darśa}) della pietra di \textit{padmarāga}; \textit{padmarāga} indica una gemma di tonalità rosso-loto, spesso intesa nella tradizione devozionale come rubino o gemma rosso-loto;}%
\wordline{नव-विद्रुम-बिम्ब-श्री-न्यक्कारि-रदन-च्छदा}{nava-vidruma-bimba-śrī-nyakkāri-radana-cchadā}
  { il cui rivestimento dei denti (\textit{radana-cchadā}, le labbra) è tale da mettere in ombra (\textit{nyakkāri}; letteralmente: “atto a far sfigurare, mettere in ombra”) la bellezza (\textit{śrī}) del corallo fresco (\textit{nava-vidruma}) e del frutto \textit{bimba}: \textit{Coccinia grandis}, pianta rampicante diffusa in India, dal frutto piccolo e tondeggiante, di colore rosso vivo a maturazione, tradizionale termine di paragone poetico per labbra fresche, piene e lucenti.}%
}

\textit{Śloka} 9 prosegue il movimento discendente della contemplazione, passando dalle orecchie alle guance e alle labbra. Le guance (\textit{kapola}) sono descritte come capaci di mettere in ombra persino la lucentezza levigata della gemma, indicando una bellezza che non si limita a riflettere, ma eccede ogni modello artificiale.

Il secondo \textit{pāda} si concentra sulle labbra, definite come “rivestimento dei denti” (\textit{radana-cchadā}). Il verbo \textit{nyakkāri} introduce un valore comparativo forte: non una semplice somiglianza, ma la capacità di far sfigurare i termini di paragone tradizionali — corallo e frutto \textit{bimba}. La bellezza di \textit{Lalitā} non si misura per confronto: essa annulla il confronto stesso.

Questo \textit{śloka} contiene due \textit{nāma}.  
\textbf{\textit{Padmarāga-śilā-darśa-\\ paribhāvi-kapola-bhūḥ}} indica che le guance superano la brillantezza della gemma.  
\textbf{\textit{Nava-vidruma-bimba-śrī-nyakkāri-radana-cchadā}} presenta le labbra come tali da mettere in ombra corallo e frutto, esprimendo una vitalità che eccede ogni canone poetico.
